% BHU PYQ -- Manually transcribed & error-corrected
% Paper: BCH-213 : Specialised Accounts-I
% Exam: B.Com. (Hons.) Semester III Examination, 2022-23

\documentclass[11pt,a4paper]{article}
\usepackage[a4paper,top=2.5cm,bottom=2.5cm,left=2.8cm,right=2.8cm,headheight=14pt]{geometry}
\usepackage{amsmath,amssymb}
\usepackage[T1]{fontenc}
\usepackage[utf8]{inputenc}
\usepackage{lmodern,microtype,fancyhdr,enumitem,hyperref,lastpage,parskip,booktabs}

\hypersetup{
    colorlinks=true,
    urlcolor=black,
    linkcolor=black,
    pdftitle={BCH-213: Specialised Accounts-I -- BHU B.Com. (Hons.) Sem III 2022-23}
}

\pagestyle{fancy}
\fancyhf{}
\renewcommand{\headrulewidth}{0.4pt}
\renewcommand{\footrulewidth}{0.4pt}
\fancyhead[L]{\small\textit{Banaras Hindu University}}
\fancyhead[R]{\small\textit{BCH-213: Specialised Accounts-I}}
\fancyfoot[C]{\small Page \thepage\ of \pageref{LastPage}}

\newlist{parts}{enumerate}{2}
\setlist[parts,1]{label=(\alph*),leftmargin=*,itemsep=5pt,topsep=3pt}
\setlist[parts,2]{label=(\roman*),leftmargin=*,itemsep=3pt,topsep=2pt}
\newcommand{\pts}[1]{\hfill\textit{[#1]}}

\begin{document}

\begin{center}
  {\large\bfseries BANARAS HINDU UNIVERSITY}\\[2pt]
  {\normalsize B.Com. (Hons.) --- Semester III Examination, 2022-23}\\[6pt]
  \rule{0.8\linewidth}{0.5pt}\\[6pt]
  {\large\bfseries Paper No. BCH-213: Specialised Accounts-I}\\[4pt]
  {\normalsize Time: 3 Hours\hfill Maximum Marks: 70}
\end{center}

\rule{\linewidth}{0.4pt}

\smallskip
\noindent\textit{Instructions: Answer all questions. All questions carry equal marks (14 marks each).}

\bigskip

\noindent\textbf{Question 1.} \\
B obtained on 1st April, 2017 from A, a lease of some coal-bearing land, the terms being a royalty of Rs. 5 a ton of coal raised subject to a minimum rent of Rs. 20,000 per annum with a right of recoupment of shortworking over the first four years of the lease.

B granted a sub-lease of part of the land to C on a royalty of Rs. 7.50 a ton merging into a minimum rent of Rs. 10,000 per annum with a right of recoupment of shortworking during the two years following the shortworking. The output for the first five years is as follows:

\begin{center}
  \begin{tabular}{ccc}
    \toprule
    \textbf{Year ended} & \textbf{Output of B (in Tons)} & \textbf{Output of C (in Tons)} \\
    \midrule
    31.03.2018 & 2,200 & 800 \\
    31.03.2019 & 2,320 & 1,080 \\
    31.03.2020 & 2,600 & 1,400 \\
    31.03.2021 & 2,800 & 1,800 \\
    31.03.2022 & 3,600 & 2,400 \\
    \bottomrule
  \end{tabular}
\end{center}

Give the necessary Ledger Accounts and Balance Sheets in the books of B and C to record these transactions, their books being closed on 31st March each year.\pts{14}

\centerline{\textbf{OR}}

How is royalty different from rent? Explain and illustrate the Journal Entries in the books of sub-lessee. Give a specimen of Nazrana Account in the books of landlord.\pts{14}

\medskip\rule{\linewidth}{0.2pt}

\noindent\textbf{Question 2.} \\
On 1st April, 2019 Saurabh Transport Company purchases on Installment Payment System one truck for Rs. 10,20,000 payable by three equal annual installments combining principal and interest, the interest being at the rate of 5\% p.a. Calculate the amount of cash price by Annuity Method and show the necessary Journal and Ledger Accounts of three years in the books of the Saurabh Transport Company. Depreciation is charged @ 10\% p.a. on reducing balance. (The present value of an annuity of one rupee for three years @ 5\% is Rs. 2.72325).\pts{14}

\centerline{\textbf{OR}}

Mention the merits and demerits of Hire-Purchase System. How is it different from Credit Purchase System? Give a specimen of Hire-Purchase Trading Account.\pts{14}

\medskip\rule{\linewidth}{0.2pt}

\noindent\textbf{Question 3.} \\
On 31st March, 2014 the following Trial Balance was sent by the New Delhi Branch to its Head Office in New York:

\begin{center}
  \small
  \begin{tabular}{lrr}
    \toprule
    \textbf{Particulars} & \textbf{Debit (Rs.)} & \textbf{Credit (Rs.)} \\
    \midrule
    Stock on 01.04.2013 & 14,400 & -- \\
    Purchases and Sales & 1,62,000 & 3,00,000 \\
    Wages and salaries & 12,000 & -- \\
    Sundry charges & 7,500 & -- \\
    Commission allowed and received & 1,500 & 4,800 \\
    Bad debts & 1,350 & -- \\
    Depreciation on furniture & 480 & -- \\
    Debtors and creditors & 19,170 & 13,500 \\
    Bills receivable and bills payable & 8,100 & 10,800 \\
    Cash in hand & 10,800 & -- \\
    Furniture & 7,200 & -- \\
    Plant (purchased on 01.10.2013) & 33,000 & -- \\
    Machinery & 1,05,000 & -- \\
    New York H.O. Account & -- & 53,400 \\
    \midrule
    \textbf{Total} & \textbf{3,82,500} & \textbf{3,82,500} \\
    \bottomrule
  \end{tabular}
\end{center}

Closing stock on 31st March, 2014 was Rs. 1,08,000. The New Delhi Branch A/c and Branch Machinery A/c appeared as \$1,500 and \$3,000 respectively in the H.O. Books.

The rates of exchange were as follows:
\begin{itemize}
  \item 01.04.2013: Rs. 40 per dollar
  \item 31.03.2014: Rs. 45 per dollar
  \item Average rate for the year 2013-14: Rs. 50 per dollar
  \item 01.10.2013: Rs. 55 per dollar
\end{itemize}
Convert the above Trial Balance and prepare the Trading and P \& L A/c and Balance Sheet of the New Delhi Branch in the New York Books.\pts{14}

\centerline{\textbf{OR}}

Point out the concept and types of Branch. How is Branch Account prepared? Explain and illustrate.\pts{14}

\medskip\rule{\linewidth}{0.2pt}

\noindent\textbf{Question 4.}
\begin{parts}
  \item The Directors of Departmental Stores Ltd. wish to ascertain the estimated net profits of the A, B and C departments separately for four months ended 31st July, 2022. It is found impracticable actually to take stock on that date but an adequate system of departmental accounting is in use and the normal rates of gross profit for the departments concerned are 40\%, 30\% and 20\% on turnover, respectively.
  
  Indirect expenses are charged in proportion to departmental turnover. For the above-mentioned period, total indirect expenses (including those relating to other departments) are Rs. 42,000 on the total turnover of Rs. 8,40,000. The following are the figures for the three departments:
  
  \begin{center}
    \begin{tabular}{lrrr}
      \toprule
      \textbf{Particulars} & \textbf{Dept. A (Rs.)} & \textbf{Dept. B (Rs.)} & \textbf{Dept. C (Rs.)} \\
      \midrule
      Stock on 1st April, 2022 & 60,000 & 70,000 & 30,000 \\
      Purchases for the period & 70,000 & 65,000 & 47,000 \\
      Sales for the period & 1,20,000 & 1,00,000 & 60,000 \\
      Direct expenses & 20,200 & 14,500 & 7,100 \\
      \bottomrule
    \end{tabular}
  \end{center}
  
  Prepare a Statement for the Directors, making a stock reserve of 10\% for each department on 31st July, 2022.\pts{10}
  
  \item State the procedure of Departmental Accounting.\pts{4}
\end{parts}

\centerline{\textbf{OR}}

Distinguish between `Department and Branch' and also between `Departmental Accounts and Branch Accounts'. State the various bases for allocation of common expenses to various departments.\pts{14}

\medskip\rule{\linewidth}{0.2pt}

\noindent\textbf{Question 5.} \\
Mention the concept, advantages and limitations of Human Resource Accounting. What are its models? Briefly explain.\pts{14}

\centerline{\textbf{OR}}

Write lucid notes on any two of the following:\pts{7+7}
\begin{parts}
  \item Methods of Human Resource Accounting
  \item Royalty accounts
  \item Independent Branch
\end{parts}

\vfill
\rule{\linewidth}{0.4pt}
\begin{center}\small
  \href{https://bhu-exam.vercel.app/}{Click here to visit our website}\\[4pt]
  \href{https://me-aryan.vercel.app/}{Click here to visit our contributor}\\[4pt]
  \href{https://quashan-me.vercel.app/}{Click here to visit our contributor}
\end{center}

\end{document}
